\section{Уводни бележки. Базов Python.}


\subsection{Уводни бележки}
	Мотивация - Курсът има широко приложение в индустрията.\\	
	Литература:
	\begin{enumerate}
	\item  Introduction to statistical learning \cite{james2023introduction}.
	\item  Elements of Statistical learning  \cite{hastie01statisticallearning}.
	\end{enumerate}
	
	\noindent
	Какво ще съдържа курса? Линейни модели - линейна и логистична регресия, Lasso, Ridge. Нелинейни модели - Splines, Обобщении адитивни модели.
	Модели с дървета - Decision trees, Random forest. Support Vector Machines. Клъстериране.   \\
	
	\noindent
	Исторически бележки - Vapnik theory.(може да се добавят няколко изречения) \\

	\noindent	
	Среда за писане на код - VsCode или Spyder, може и Jupyter Notebook. Anaconda не се препоръчва заради лицензи. Линк за повече инфо. \\
	Инсталация с uv? \\
 	Тук ще изпозлваме Python 3.10.11, последната версия с prebuild инсталационен пакет от python.org,
 	от версиите 3.10.11. За 3.11 и наскоро излязлата 3.12 нямаме още всички пакети. Ако ще изпозлваме виртуална среда: 	
 	\begin{lstlisting}[language=bash]
$pip install virtualenv
 	\end{lstlisting}
 	(Да се добавят git clone и прочие... команди и и описания) \\
 	
 	Отваряме с VSCode папката с кода на упражненията и пишем:
	\begin{lstlisting}[language=bash]
$pip install -r requirements.txt
 	\end{lstlisting}
 	За инсталиране на пакетите, които се използват в курса. \\
 	Линк към страницата с Lab упражненията от книгата:\\
 	\href{https://intro-stat-learning.github.io/ISLP/installation.html}{https://intro-stat-learning.github.io/ISLP/installation.html}
 	
 	\subsection{Базов Python}
 	Най-базови неща в Python - python list, set, import packages, най-базови фунцкии в numpy,
 	събиране на вектори и умножение на вектори с число, скаларно произведение, матрици.
 	pandas пакета и пример за графика с matplotlib.pyplot.\\
 	
 	Lab 2.3 Introduction to Python от Python книгата. Линк: \\

 	\href{https://intro-stat-learning.github.io/ISLP/labs/Ch02-statlearn-lab.html#}{https://intro-stat-learning.github.io/ISLP/labs/Ch02-statlearn-lab.html\#}.